\subsection{General Structure of a VQA}\label{sec:general-structure-of-a-vqa}

\subsubsection{Variational Quantum Algorithms}\label{sec:variational-quantum-algorithms}

A \definition{Variational Quantum Algorithm (VQA)} is a hybrid algorithm in which:
\begin{enumerate}
    \item A \textbf{parameterized quantum circuit} generataes candidate solutions to a problem;
    \item A \textbf{classical optimizer} evaluates those solutions;
    \item The process is repeated until the best parametes are found, which correspond to the optimal solution of the problem.
\end{enumerate}
The key idea is that we do \textbf{not directly program the solution}. Instead, we program a \hl{circuit whose behavior depends on adjustable parameters}.

\highspace
\begin{flushleft}
    \textcolor{Green3}{\faIcon{sitemap} \textbf{The Three Components}}
\end{flushleft}
Every VQA contains exactly three fundamental ingredients:
\begin{itemize}
    \item \important{Parametric Quantum Circuit (Ansatz)}. The \hl{quantum computer executes a circuit containing adjustable parameters}. For example, if the circuit contains a rotation gate $R_y(\theta)$, then $\theta$ is a parameter that can be tuned to change the behavior of the circuit. The quantum computer therefore acts as a \textbf{state generator}, which produces different quantum states depending on the values of the parameters. We will see the ansatz in detail in the next section.


    \item \important{Classical Optimizer}. The optimizer \textbf{chooses the values of the parameters}. It is a sort of coach that guides the quantum computer towards the optimal solution. The optimizer asks itself: ``\emph{which $\theta$ should I try next?}''. As we have seen, the optimizer runs on a classical computer, and it can be implemented using various algorithms, such as gradient descent, genetic algorithms, or Bayesian optimization.


    \item \important{Iterative Execution}. The VQA is an \textbf{iterative process}: choose parameters, run the quantum circuit, measure the output, evaluate the cost function, update the parameters, and repeat if necessary. This loop continues until some stopping condition is reached.
\end{itemize}

\highspace
\begin{flushleft}
    \textcolor{Green3}{\faIcon{question-circle} \textbf{Why \emph{Variational}?}}
\end{flushleft}
The name \emph{variational} comes from the fact that the \textbf{parameters vary}. During the execution of the algorithm, we are constantly changing the parameters of the quantum circuit in order to find the optimal solution. Thus, the term \emph{variational} emphasizes the \textbf{algorithm's dynamic nature}. Rather than running a fixed circuit, the \hl{algorithm searches through a family of quantum states generated by different parameter values}.

\highspace
In summary, \definition{Variational Quantum Algorithms (VQAs)} are \textbf{hybrid\break quantum-classical algorithms} consisting of three main components: a parameterized quantum circuit called an Ansatz, a classical optimizer that searches for the best parameters, and an iterative execution loop. The quantum computer evaluates candidate solutions, while the classical optimizer updates the circuit parameters according to the measured results until an optimal solution is found.